\hypertarget{discussion}{%
\section{Discussion}}

A strong size-based forecast changes the interpretation of the timing gain. Conditioning the empirical distribution on prefix size reduces CRPS by approximately 18\% on Weibo and 24\% on SEISMIC. QRF with timing improves further on that conditional reference by 6.91\% and 8.61\%, respectively. These descriptive procedure comparisons show why an unconditional reference alone can make a richer learner look more advantageous. They do not partition predictive information or identify the fraction of achievable accuracy already exhausted.

Within each primary family, adding timing improves average count CRPS on both sources. The magnitude depends on the source and estimator, and the supplementary strata show that this average conceals mixed subgroup responses. Dense-prefix QRF improvements recur across the two sources, while other families and sparse prefixes do not collapse to the same pattern. The contrast at empty prefixes further shows that changing a model's training representation can change predictions where the additional local measurements are absent. The completed control strengthens the aggregate finding: the specified timing representation outperforms equally many independent noise coordinates in every source--family case. This result and the empty-prefix counterexample address different questions. The former supports the predictive usefulness of the trained timing representation; the latter shows why a subgroup contrast cannot automatically be attributed to event-time information observed locally.

One possible mechanism in Hurdle's linear predictors is that the sentinel-coded timing features provide additional directions for separating empty from nonempty prefixes, permitting a regime-like response beyond a linear term in log(1+n). Emptiness is already recoverable from I0; this hypothesis concerns representation and fitted flexibility rather than new information, and its role in the observed improvement was not isolated experimentally.

This interpretation connects timing-feature research with the distinction between predictive performance and intrinsic predictability. The result is an estimate of what the specified procedures gain from the representation change on these samples. Remaining error could reflect model restrictions, training data, archive selection or irreducible variation \cite{martin2016,hofman2017}.

Calibration gives a further practical qualification. Lower CRPS and improved no-growth loss can coexist with uneven interval coverage and persistent PIT departures. Comparing scores and calibration together distinguishes average distributional improvement from calibrated probabilities in every subgroup.

The sources provide different observation contexts: one nonroot event is the median five-minute prefix on Weibo, compared with 30 on SEISMIC. Stratification makes these differences visible without isolating a platform mechanism. Selection criteria, response scales, split design and model fitting remain intertwined. Full-history point processes and specialized cascade architectures address a broader information contract; the present study offers a restricted-input empirical reference for those comparisons. {\small\textit{Preparation:~\cite{openai2026}.}}
